\subsection{\texorpdfstring{HAUSDORFF QUOTIENT GROUPS OF \(R^{n}\)}{HAUSDORFF QUOTIENT GROUPS OF R TO THE N}}

Every Hausdorff quotient group of \(\mathbf{R}^n\) is of the form \(\mathbf{R}^n/\mathbf{H}\) , where \(\mathbf{H}\) is a closed subgroup of \(\mathbf{R}^n\) (Chapter III, § 2, no. 6, Proposition 18). By Corollary 2, Theorem 2 of no. 3, there exists an automorphism \(f\) of \(\mathbf{R}^n\) which transforms \(\mathbf{H}\) into a subgroup \(\mathbf{H}'\) , the direct sum of a vector subspace generated by \(p\) of the vectors \(\boldsymbol{e}_i\) of the canonical basis and the discrete group generated by \(q\) of the \(n-p\) remaining vectors

\[
\boldsymbol {e} _ {i}, \quad \mathrm{o} \leqslant p + q \leqslant n.
\]

Passing to the quotients, \(f\) induces an isomorphism of \(\mathbf{R}^n / \mathbf{H}\) onto \(\mathbf{R}^n / \mathbf{H}'\) (Chapter III, § 2, no. 8, Remark 3); now, \(\mathbf{R}^n / \mathbf{H}'\) is isomorphic to \(\mathbf{R}^{n - p - q} \times \mathbf{T}^q\) (Chapter III, § 2, no. 9, Corollary to Proposition 26). Consequently:

PROPOSITION 9. Every Hausdorff quotient group of \(\mathbf{R}^{n}\) is isomorphic to a product group \(\mathbf{R}^{h} \times \mathbf{T}^{k}\) ( \(0 \leqslant h + k \leqslant n\) ).

The product space \(\mathbf{T}^n\) (and, by abuse of language, the topological group \(\mathbf{T}^n\) ) is called the \(n\) -dimensional torus; by Proposition 4 of Chapter V, §1, no. 2, \(\mathbf{T}^n\) is compact, connected and locally connected.

Furthermore, if C denotes a closed cube of side 1 in \(R^{n}\) , \(T^{n}\) is homeomorphic to the quotient space of C by the equivalence relation `` \(x_{i} \equiv y_{i} \pmod{1}\) ( \(1 \leqslant i \leqslant n\) )'' between the points \(\boldsymbol{x} = (\boldsymbol{x}_{i})\) and \(\boldsymbol{y} = (\boldsymbol{y}_{i})\) of C. Thus \(T^{n}\) is formed from the cube C by ``identification of opposite faces''.

PROPOSITION 10. The topological group \(\mathbf{T}^n\) is locally isomorphic to \(\mathbf{R}^n\) . For \(\mathbf{T}^n = (\mathbf{R}/\mathbf{Z})^n\) is isomorphic to \(\mathbf{R}^n/\mathbf{Z}^n\) (Chapter III, § 2, no. 9, Corollary to Proposition 26) and \(\mathbf{Z}^n\) is a discrete subgroup of \(\mathbf{R}^n\) (Chapter III, § 2, no. 6, Proposition 19).

It follows that the groups \(\mathbf{R}^p\times \mathbf{T}^{n - p}\) are all locally isomorphic to \(\mathbf{R}^n\) \((\mathrm{o}\leqslant p\leqslant n)\) ; in § 2, no. 2, we shall see that they are the only connected groups which have this property.

